\documentclass[
  manuscript=article,  %% article
  layout=preprint,  %% preprint (for submission) or publish (for publisher only and after acceptance)
  year=20xx,
  volume=x,
  IssueNumber=x
]{extra/mjcis}

\doi{xx.xxxxx/MJCIS.xxxx.xxxx}

% only used by journal admin
\StartPage{1}



\received {1 April 20xx}
\accepted {10 May 20xx}
\published{20 May 20xx}



% --- blew is the area for authors ---

% remove the following two packages, and delete all \blindtext commands
\usepackage[english]{babel} 
\usepackage{blindtext}


% specify the .bib file for references
\addbibresource{reference.bib} 


% Make sure your article tile is within 12 words
\title{Using a concise title for your article}

\author{First Author}
\affiliation{Institution-1, City, Country}
\email{correspondence@email.domain}

\author{Second Author}
\affiliation{Institution-2, City, Country}

\author{Third Author}
\alsoaffiliation{Institution-1, City, Country}
\affiliation{Institution-3, City, Country}


% maximum five keywords
\keywords{keyword; keyword-two; keyword number three} 



\begin{document}

\begin{abstract}
  An abstract summarizes in one paragraph with 300 words or less, the major aspects of the entire paper. They often include: 1) the overall purpose of the study and the research problem you investigated; 2) the basic design of you research approach; 3) major findings as a result of your analysis; and, 4) a brief summary of your interpretations and conclusions. 
\end{abstract}


\section{Introduction}

\blindtext 
open data in science \cite{murray2008open}. 


\blindtext [2]


\section{Method}

\subsection{Method part 1}

\blindtext The end result is in Figure \ref{fig:logo}.

\begin{figure}[ht!]
  \centering
  \includegraphics[width=0.45\textwidth]{example-image}
  \caption{An Example Figure}
  \label{fig:logo}
\end{figure}

\blindtext\footnote{This is how a footnote works.}

\subsection{Method part 2}

\subsubsection{Method part 2-1}

\blindtext Reference to Equation \ref{eq:cauchy_momentum}.

\begin{equation} \label{eq:cauchy_momentum}
\rho\frac{\mathrm{D} \mathbf{u}}{\mathrm{D} t} = - \nabla p + \nabla \cdot \boldsymbol \tau + \rho\,\mathbf{g}
\end{equation}


\subsubsection{Method part 2-2}

\blindtext Table \ref{tb:example_table} shows an example.

\begin{table}[H]
  \centering
  \small
  \caption{Example table}
  \label{tb:example_table}
  \begin{tabular}{lll}
  \toprule
  \textbf{Parameter} & \textbf{Notation} & \textbf{Remarks} \\
  \midrule
  name & - & engine common identifier \\
  manufacture & - & name of the manufacture  \\
  bpr & $\lambda$ & bypass ratio \\
  pr & - & pressure ratio \\
  thrust & $T_0$ & maximum static thrust\\
  \bottomrule
  \end{tabular}
\end{table}

\blindtext


\section{Discussions}

\paragraph{Paragraph title} This is the paragraph with title if you want to use such function in the paper. \blindtext


\section{Conclusion}

\blindtext


%\appendix
%\section{Supplementary figures}
%\blindtext

%\section{Supplementary tables}
%\blindtext


\section*{Acknowledgement}
Include your acknowledgement in this section.


\section*{Funding statement}
When applicable, please specify the funding information for this research.


% Data statement is mandatory for all papers
\section*{Open data statement}
short description to supplementary data.





\printbibliography



\end{document}